% ============================================================================
% Research paper template (see 00-RESEARCH-DOCUMENTATION.md, "From study to
% paper"). Copy to <repo>/docs/research/papers/NN-SLUG.tex and replace every
% ALL-CAPS placeholder.
%
% THE HOUSE STYLE IS NOT IN THIS FILE. It is in HOUSE-PREAMBLE.tex, which you
% copy into the repository once as docs/research/papers/house-preamble.tex and
% every document there \inputs -- so that fixing the house style is one edit
% rather than one edit per document, and so that a document cannot drift out
% of it by being written from memory. That file's header states the contract;
% the four \newcommands below are it.
%
% Build (PDFs are build output -- never committed):
%   pdflatex NN-SLUG.tex && pdflatex NN-SLUG.tex
% ============================================================================
\documentclass[11pt,a4paper]{article}

% ----------------------------------------------------------------------------
% The house preamble's contract.
%
% \paperstatus is this standard's own marker: draft | verified, mirroring the
% companion study's status header. WHILE THE STUDY IS NOT VERIFIED THIS MUST
% SAY DRAFT -- it rides in the header of every page, so a reader who opens the
% PDF at page 9 still sees it.
% ----------------------------------------------------------------------------
\newcommand{\paperstatus}{DRAFT}
\newcommand{\headerleft}{PROJECT-NAME --- SHORT-PAPER-TITLE}
\newcommand{\pdftitleline}{PAPER-TITLE}
\newcommand{\pdfauthorline}{AUTHOR-NAME}

% A second language edition of this paper adds two lines here -- and nothing
% else -- before the \input (see the standard, "Publishing a document in more
% than one language"):
%   \def\houselang{polish}
%   \def\editionsuffix{.pl}

% Packages beyond the house set go AFTER this line, with a comment saying so.
\input{house-preamble}

\begin{document}

% ============================================================================
% TITLE BLOCK. The fifth argument is the provenance line: a paper is a
% presentation of a study, and this is where it says which one and at what
% commit, so that a reader who finds the two disagreeing knows the study is
% the document and the paper is stale.
% ============================================================================
\housetitle
    {PAPER-TITLE}
    {OPTIONAL-SUBTITLE-OR-LEAVE-EMPTY}
    {AUTHOR-NAME}
    {AFFILIATION-OR-REPO-URL}
    {companion study: \texttt{docs/research/NN-SLUG.md} at commit \texttt{COMMIT}}

% ----------------------------------------------------------------------------
% Abstract: question, method, headline result. Written last, placed first --
% same rule as the markdown study.
% ----------------------------------------------------------------------------
\begin{abstract}
\noindent ABSTRACT-TEXT.
\end{abstract}

\vspace{0.4cm}

\section{Introduction}
% The context and the research questions, numbered as in the study (RQ1,
% RQ2, ...) so Results can answer them by name.

INTRODUCTION-TEXT.

\begin{itemize}[leftmargin=*]
    \item \textbf{RQ1} --- FIRST-QUESTION.
    \item \textbf{RQ2} --- SECOND-QUESTION.
\end{itemize}

\section{Background and related work}
% What was already known: the repo's theory docs, prior studies, external
% literature \cite{ref1}. Cite, don't re-teach.

BACKGROUND-TEXT.

\section{Method}
% Precise enough for a skeptic to attack. Include how the instrument was
% validated (the standard's rule 3) before its readings counted as evidence.

METHOD-TEXT.

% A figure is called by slug, never by path -- \dgm resolves the rendered file
% for this edition. Run the repository's render step BEFORE pdflatex, or the
% build dies here rather than on the missing diagram you would rather hear
% about.
%
% \begin{figure}[ht]
% \centering
% \includegraphics[width=0.7\linewidth]{\dgm{SLUG}}
% \caption{CAPTION-TEXT.}
% \end{figure}

\section{Results}
% Findings tied to RQs, evidence inline. Every number in this section is
% traceable to the companion study's tests, artifacts, or shown derivations
% -- a paper introduces no numbers of its own.

RESULTS-TEXT.

\begin{table}[h]
\centering
\begin{tabular}{@{}lrr@{}}
\toprule
\textbf{Quantity} & \textbf{Value} & \textbf{Evidence} \\
\midrule
EXAMPLE-ROW & 0.000000 & artifact / test / derivation \\
\bottomrule
\end{tabular}
\caption{CAPTION-TEXT.}
\end{table}

\section{Discussion}

\subsection{Threats to validity}
% Stated by the author before a reviewer has to.

THREATS-TEXT.

\subsection{Implications}
% What should change because of this.

IMPLICATIONS-TEXT.

\section{Conclusion}

CONCLUSION-TEXT.

\section*{Reproducibility}
% Verbatim commands -- same as the study's Reproduction section.

Every number in this paper regenerates from the companion study's committed
artifacts:

\begin{quote}
\repopath{REPRODUCTION-COMMAND}
\end{quote}

% ============================================================================
% References -- numbered and resolvable, mirroring the study's list.
% thebibliography (no BibTeX) keeps the paper a single self-contained file.
% ============================================================================
\begin{thebibliography}{9}

\bibitem{ref1}
A.~Author --- \textit{Title}, Publisher, YEAR.

\bibitem{ref2}
REPO-NAME repository --- \texttt{docs/research/NN-SLUG.md} and its
artifacts, \url{https://github.com/OWNER/REPO}.

\end{thebibliography}

\end{document}
